%% kapa: exam=2 round=34 year=2023 period=3 subject=compensation_law subject_ko=감정평가 및 보상법규 seq=1 number=1 points=40
%% source: Q-Net 감정평가사 기출문제 (cst00309) artlSeq=5212168 file=raw/1008/5212168/3교시 23년도 감정평가사 2차시험_감정평가 및 보상법규.pdf
%% conversion: kordoc 4.24.51 -> tools/gen2.py (auto); review=proofread-gpt-sol; flags=none
\begin{longquestion}{【문제 1】}
A대도시의 시장은 ｢국토의 계획 및 이용에 관한 법률｣에 따른 도시관리계획으로 관할구역 내 ○○동 일대 90,000 m² 토지에 공영주차장과 자동차정류장을 설치하는 도시계획시설사업결정을 한 후 ｢지방공기업법｣에 따른 A대도시 X지방공사(이하 ‘X공사’라 함)를 도시계획시설사업의 시행자로 지정하고, X공사가 작성한 실시계획에 대해 실시계획인가를 하고 이를 고시하였다. 이에 따라 ｢공익사업을 위한 토지 등의 취득 및 보상에 관한 법률｣(이하 ‘토지보상법’이라 함)에 의해 사업인정 및 고시가 이루어졌다. 한편, X공사는 사업대상구역 내에 위치한 20,000 m² 토지를 소유한 甲과 토지수용을 위한 협의를 진행하였으나 협의가 성립되지 아니하여 관할 지방토지수용위원회에 토지수용의 재결을 신청하였다. 다음 물음에 답하시오. (단, 각 물음은 상호독립적임) (40점) \par
\medskip\noindent 물음 1) 토지보상법상의 사업인정과 사업인정고시의 법적 성질에 관하여 설명하시오. (10점) \par
\medskip\noindent 물음 2) 甲은 수용 자체가 위법이라고 주장하면서 관할 지방토지수용위원회의 수용재결과 중앙토지수용위원회의 이의재결을 거친 후 취소소송을 제기하였다. 취소소송의 대상적격과 피고적격에 관하여 설명하시오. (20점) \par
\medskip\noindent 물음 3) 甲은 자신의 토지에 대한 보상금이 적으며, 일부 지장물이 손실보상 대상에서 제외되었다는 이유로 관할 지방토지수용위원회의 수용재결에 불복하여 중앙토지수용위원회에 이의신청을 거쳤으나 기각재결을 받았다. 甲이 이에 대하여 불복하는 경우 적합한 소송 형태를 쓰고 이에 관하여 설명하시오. (10점) \par
\begin{framedtext}
참조 조문 \par
<국토의 계획 및 이용에 관한 법률> \par
제88조(실시계획의 작성 및 인가 등) ① 도시ㆍ군계획시설사업의 시행자는 대통령령으로 정하는 바에 따라 그 도시ㆍ군계획시설사업에 관한 실시계획(이하 “실시계획”이라 한다)을 작성하여야 한다. \par
② 도시ㆍ군계획시설사업의 시행자(국토교통부장관, 시ㆍ도지사와 대도시 시장은 제외한다. 이하 제3항에서 같다)는 제1항에 따라 실시계획을 작성하면 대통령령으로 정하는 바에 따라 국토교통부장관, 시ㆍ도지사 또는 대도시 시장의 인가를 받아야 한다. 다만, 제98조에 따른 준공검사를 받은 후에 해당 도시ㆍ군계획시설사업에 대하여 국토교통부령으로 정하는 경미한 사항을 변경하기 위하여 실시계획을 작성하는 경우에는 국토교통부장관, 시ㆍ도지사 또는 대도시 시장의 인가를 받지 아니한다. \par
제96조(「공익사업을 위한 토지 등의 취득 및 보상에 관한 법률」의 준용) ① 제95조에 따른 수용 및 사용에 관하여는 이 법에 특별한 규정이 있는 경우 외에는 「공익사업을 위한 토지 등의 취득 및 보상에 관한 법률」을 준용한다. \par
② 제1항에 따라 「공익사업을 위한 토지 등의 취득 및 보상에 관한 법률」을 준용할 때에 제91조에 따른 실시계획을 고시한 경우에는 같은 법 제20조제1항과 제22조에 따른 사업인정 및 그 고시가 있었던 것으로 본다. 다만, 재결 신청은 같은 법 제23조제1항과 제28조제1항에도 불구하고 실시계획에서 정한 도시ㆍ군계획시설사업의 시행기간에 하여야 한다. \par
<공익사업을 위한 토지 등의 취득 및 보상에 관한 법률> \par
제28조(재결의 신청) ① 제26조에 따른 협의가 성립되지 아니하거나 협의를 할 수 없을 때(제26조제2항 단서에 따른 협의 요구가 없을 때를 포함한다)에는 사업시행자는 사업인정고시가 된 날부터 1년 이내에 대통령령으로 정하는 바에 따라 관할 토지수용위원회에 재결을 신청할 수 있다. \par
제83조(이의의 신청) ① 중앙토지수용위원회의 제34조에 따른 재결에 이의가 있는 자는 중앙토지수용위원회에 이의를 신청할 수 있다. \par
② 지방토지수용위원회의 제34조에 따른 재결에 이의가 있는 자는 해당 지방토지수용위원회를 거쳐 중앙토지수용위원회에 이의를 신청할 수 있다. \par
③ 제1항 및 제2항에 따른 이의의 신청은 재결서의 정본을 받은 날부터 30일 이내에 하여야 한다. \par
제84조(이의신청에 대한 재결) ① 중앙토지수용위원회는 제83조에 따른 이의신청을 받은 경우 제34조에 따른 재결이 위법하거나 부당하다고 인정할 때에는 그 재결의 전부 또는 일부를 취소하거나 보상액을 변경할 수 있다. \par
② 제1항에 따라 보상금이 늘어난 경우 사업시행자는 재결의 취소 또는 변경의 재결서 정본을 받은 날부터 30일 이내에 보상금을 받을 자에게 그 늘어난 보상금을 지급하여야 한다. 다만, 제40조제2항 \par
제1호ㆍ제2호 또는 제4호에 해당할 때에는 그 금액을 공탁할 수 있다. \par
\end{framedtext}
\end{longquestion}
